\section{Branch-and-Bound}
Branch-and-Bound (B\&B) is the foundation of this project and is one of the numerous techniques available for solving combinatorial optimization problems, such as the scheduling problem. Other techniques include greedy algorithms, dynamic programming, and heuristic methods, which we will not mention in this context.

The Branch-and-Bound method is used to solve NP-hard problems. The algorithm systematically explores the solution space by dividing the problem into smaller subproblems (\emph{branching}) and calculating lower and upper bounds to evaluate the quality of the solutions found (\emph{bounding}). Once a subproblem's bound indicates that it cannot yield a better solution than the best one found so far, a new subproblem is selected for exploration. If a subproblem's bound is worse than the current best solution, it is ignored and the process continues (\emph{selection/search}). This process continues until all subproblems have been explored or pruned, ultimately leading to the optimal solution.

From the brief description above, we can identify three main components of Branch-and-Bound:
\begin{itemize}
  \item \textbf{Branching}: The division of the problem into smaller subproblems. In our context, this means to assign a specific job to a specific machine. We will refer to this as \emph{fixing} a job to a machine.
  \item \textbf{Bounding}: The calculation of lower and upper bounds for each subproblem. These bounds help determine whether a subproblem can be discarded (pruning) or if it should be further explored.
  \item \textbf{Selection/Search}: The strategy used to select which subproblem to explore next. This can be done using different strategies, such as depth-first or breadth-first.
\end{itemize}

To better understand the algorithm's flow and its application to the scheduling problem, we delve into its specific characteristics in the next section.


\subsection{Mathematical Formulation}
To precisely describe the Branch-and-Bound algorithm for the identical machines scheduling problem, we introduce the following mathematical notation:
\begin{itemize}
    \item $x_{j,i} \in [0,1]$ indicates the fraction of job $j$ assigned to machine $i$
    \item $F \subseteq \{1,\ldots,n\} \times \{1,\ldots,m\}$ is the set of fixed assignments at a given node
    \item For a node with fixed assignments $F$, the residual problem considers only the unfixed jobs
    \item $LB(F)$ is the lower bound on the makespan for the subproblem with fixed assignments $F$
    \item $UB(F)$ is the upper bound on the makespan obtained by rounding a fractional solution
    \item $\text{GLB}, \text{GUB}$ are the global lower and upper bounds on the optimal makespan
\end{itemize}

The algorithm maintains the following invariant: at any point during execution, the optimal makespan $C^*$ satisfies:
\begin{equation}
\text{GLB} \leq C^* \leq \text{GUB}
\end{equation}

\begin{algorithm}[H]
\KwIn{Instance with $n$ jobs, processing times $p_1, \ldots, p_n$, and $m$ identical machines}
\KwOut{Job-machine assignment minimizing makespan with quality guarantee $(1+\varepsilon)$}

\tikzmk{A}\algcomment{Initialization}\tikzmk{B}\boxit{mygray}\\
Solve relaxed problem for root node $F = \emptyset$ to obtain lower bound $LB(\emptyset)$\;
Round fractional solution to obtain feasible assignment with upper bound $UB(\emptyset)$\;
Initialize queue $\mathcal{Q} \leftarrow \{(\emptyset, LB(\emptyset))\}$\;
Set global bounds: $\text{GLB} \leftarrow LB(\emptyset)$, $\text{GUB} \leftarrow UB(\emptyset)$\;
\BlankLine
\While{$\mathcal{Q} \neq \emptyset$ and $\frac{\text{GUB}}{\text{GLB}} \geq 1 + \varepsilon$}{
    \tikzmk{A}\algcomment{Selection}\tikzmk{B}\boxit{mygray}\\
    Select node $(F, LB(F))$ from $\mathcal{Q}$ according to search strategy\;
    
    \BlankLine
    \tikzmk{A}\algcomment{Branching}\tikzmk{B}\boxit{mygray}\\
    Choose fractional job $j^* \in \{1,\ldots,n\} \setminus \{j : (j,i) \in F \text{ for some } i\}$\;
    
    \BlankLine
    \tikzmk{A}\algcomment{Bounding and Pruning}\tikzmk{B}\boxit{mygray}\\
    \ForEach{machine $i \in \{1,\ldots,m\}$}{
        Construct child node $F' \leftarrow F \cup \{(j^*, i)\}$\;
        
        \eIf{$|F'| = n$}{
            Compute makespan $C(F')$ of complete assignment $F'$\;
            \If{$C(F') < \text{GUB}$}{
                $\text{GUB} \leftarrow C(F')$\;
            }
        }{
            Solve relaxed problem for $F'$ to obtain $LB(F')$\;
            \If{$LB(F') < \text{GUB}$}{
                Round to obtain feasible solution with $UB(F')$\;
                $\text{GUB} \leftarrow \min\{\text{GUB}, UB(F')\}$\;
                Add $(F', LB(F'))$ to $\mathcal{Q}$\;
            }
        }
    }
    
    \BlankLine
    \tikzmk{A}\algcomment{Update global lower bound}\tikzmk{B}\boxit{mygray}\\
    $\text{GLB} \leftarrow \min\{\text{GUB}, \min_{(F,LB(F)) \in \mathcal{Q}} LB(F)\}$\;
}
\BlankLine
\Return{Best assignment found with makespan $\text{GUB}$}\;
\caption{Branch-and-Bound for Identical Machines Scheduling}
\end{algorithm}


\subsection{Key Mathematical Components}

\subsubsection{Lower Bound Computation}
For a node with fixed assignments $F$, the lower bound $LB(F)$ is computed by solving a relaxed version of the scheduling problem where jobs can be fractionally split across machines. Let $o_i = \sum_{(j,i) \in F} p_j$ denote the overhead (total processing time already fixed) on machine $i$. The relaxed problem seeks to minimize:
\begin{equation}
\min_{x_{j,i} \geq 0} \max_{i=1,\ldots,m} \left\{ o_i + \sum_{j \notin \{j' : (j',i') \in F\}} p_j \cdot x_{j,i} \right\}
\end{equation}
subject to:
\begin{equation}
\sum_{i=1}^m x_{j,i} = 1 \quad \forall j \notin \{j' : (j',i') \in F\}
\end{equation}

This can be efficiently solved using greedy allocation or binary search techniques.

\subsubsection{Upper Bound via Rounding}
Given a fractional solution $\{x_{j,i}\}$ from the relaxed problem, an upper bound is obtained by converting it to a complete assignment. Since at most $m$ jobs can be fractionally assigned, these are rounded to integral assignments using a matching-based approach that minimizes the resulting makespan.

\subsubsection{Pruning Criterion}
A node $(F, LB(F))$ is pruned if:
\begin{equation}
LB(F) \geq \text{GUB}
\end{equation}
This ensures that the subtree rooted at this node cannot contain a better solution than the current best.

\subsubsection{Termination Condition}
The algorithm terminates when the approximation guarantee is satisfied:
\begin{equation}
\frac{\text{GUB}}{\text{GLB}} < 1 + \varepsilon
\end{equation}
At termination, the solution with makespan $\text{GUB}$ is guaranteed to be within a factor of $(1+\varepsilon)$ of the optimal makespan $C^*$.